
\documentclass[11pt]{article}
\usepackage[ngerman]{babel}

\usepackage[a4paper,margin=0.9in]{geometry}
\usepackage[T1]{fontenc}
\usepackage[utf8]{inputenc}
\usepackage{lmodern}
\usepackage{amsmath,amssymb,mathtools}
\usepackage{microtype}
\usepackage{fancyhdr}
\usepackage{array,enumitem}
\setlength{\parindent}{1.4em}
\setlength{\parskip}{0.15em}
\emergencystretch=6em
\setlength{\headheight}{14pt}
\pagestyle{fancy}
\fancyhf{}
\cfoot{\thepage}
\newcommand{\sect}[2]{\section*{\S\,#1. #2}}
\newcommand{\ord}{\operatorname{ord}}

\lhead{Weber, Lehrbuch der Algebra I}
\rhead{Deutsche Quelle: \S\,46 Ende--\S\,54}
\begin{document}
\begin{center}
{\large Heinrich Weber}\\[0.35em]
{\Large Lehrbuch der Algebra. Erster Band.}\\[0.35em]
{\large Vierter Abschnitt. Symmetrische Functionen.}\\[0.25em]
{\large \S\,46 Ende--\S\,54}
\end{center}

\sect{46}{Discriminanten}

Hierin ist nach \S\,42 zu setzen
\[
\begin{gathered}
 s_0=3,\qquad a_0s_1=-a_1,\\
 a_0^2s_2=a_1^2-2a_0a_2,\\
 a_0^3s_3=-a_1^3+3a_0a_1a_2-3a_0^2a_3,\\
 a_0^4s_4=a_1^4-4a_0a_1^2a_2+4a_0^2a_1a_3+2a_0^2a_2^2,
\end{gathered}
\]
wodurch man erhält
\begin{equation}
D=a_1^2a_2^2+18a_0a_1a_2a_3-4a_0a_2^3-4a_1^3a_3-27a_0^2a_3^2.
\tag{10}
\end{equation}

\sect{47}{Discriminanten der Formen dritter und vierter Ordnung}

Bei den Functionen dritter und vierter Ordnung haben wir zur Berechnung der Discriminanten in den Auflösungen der Gleichungen dieser beiden Grade (\S\,35, 37) ein einfaches Mittel.

Was zunächst die Gleichung dritter Ordnung betrifft, so waren die drei Wurzeln von
\begin{equation}
x^3+ax+b=0
\tag{1}
\end{equation}
in der Form dargestellt:
\begin{equation}
\begin{aligned}
\alpha&=u+v,\\
\beta&=\varepsilon u+\varepsilon^2v,\\
\gamma&=\varepsilon^2u+\varepsilon v,
\end{aligned}
\tag{2}
\end{equation}
worin
\begin{equation}
 \varepsilon=\frac{-1+\sqrt{-3}}{2},\qquad
 \varepsilon^2=\frac{-1-\sqrt{-3}}{2}
\tag{3}
\end{equation}
die complexen dritten Einheitswurzeln sind und $u,v$ durch \S\,35, (8) bestimmt sind. Aus (3) folgt
\begin{equation}
 \varepsilon-\varepsilon^2=\sqrt{-3},\qquad
 1-\varepsilon=\varepsilon^2\sqrt{-3},\qquad
 1-\varepsilon^2=-\varepsilon\sqrt{-3},
\tag{4}
\end{equation}
und aus (2)
\[
\begin{aligned}
\alpha-\beta&=\sqrt{-3}\, (\varepsilon^2u-\varepsilon v),\\
\alpha-\gamma&=-\sqrt{-3}\,(\varepsilon u-\varepsilon^2v),\\
\beta-\gamma&=\sqrt{-3}\,(u-v).
\end{aligned}
\]
Hieraus durch Multiplication
\begin{equation}
(\alpha-\beta)(\alpha-\gamma)(\beta-\gamma)=3\sqrt{-3}\,(u^3-v^3).
\tag{5}
\end{equation}
Erhebt man ins Quadrat, so erhält man für die Discriminante der cubischen Gleichung (1) nach \S\,35, (6)
\begin{equation}
D=-(27b^2+4a^3).
\tag{6}
\end{equation}
Hieraus erhält man die Discriminante der allgemeinen Function
\begin{equation}
a_0x^3+a_1x^2+a_2x+a_3,
\tag{7}
\end{equation}
wenn man für $a$ und $b$ die Ausdrücke \S\,34, (4) einsetzt, in denen man $a_1,a_2,a_3$ durch
\[
\frac{a_1}{a_0},\quad \frac{a_2}{a_0},\quad \frac{a_3}{a_0}
\]
ersetzt und [\S\,46, (3)] mit $a_0^4$ multiplicirt. Man erhält so
\[
 a=\frac{3a_0a_2-a_1^2}{3a_0^2},
 \qquad
 b=\frac{2a_1^3-9a_0a_1a_2+27a_0^2a_3}{27a_0^3},
\]
und folglich
\begin{equation}
D=-\frac{1}{27a_0^2}\left\{(2a_1^3-9a_0a_1a_2+27a_0^2a_3)^2
       +4(3a_0a_2-a_1^2)^3\right\}.
\tag{8}
\end{equation}
Wenn man das Quadrat und den Cubus ausrechnet, so kommt man zu der Formel (10) des vorigen Paragraphen zurück; es ist aber bemerkenswerth, dass die Formel (8) scheinbar den Nenner $27a_0^2$ enthält, der sich beim Ausrechnen forthebt.

Ebenso können wir bei den biquadratischen Gleichungen verfahren. Wir erhalten, wenn $\alpha,\beta,\gamma,\delta$ die Wurzeln der biquadratischen Gleichung
\begin{equation}
x^4+ax^2+bx+c=0
\tag{9}
\end{equation}
sind, nach \S\,37, (5)
\begin{equation}
\begin{alignedat}{2}
\alpha-\beta&=v+w,          &\qquad \gamma-\delta&=v-w,\\
\alpha-\gamma&=w+u,         &\qquad \delta-\beta&=w-u,\\
\alpha-\delta&=u+v,         &\qquad \beta-\gamma&=-u+v.
\end{alignedat}
\tag{10}
\end{equation}
Danach wird die Discriminante der Gleichung (9)
\[
D=(v^2-w^2)^2(w^2-u^2)^2(u^2-v^2)^2,
\]
d. h. es ist $D$ gleich der Discriminante der cubischen Resolvente \S\,37, (4)
\begin{equation}
z^3+2az^2+(a^2-4c)z-b^2=0,
\tag{11}
\end{equation}
deren Wurzeln $u^2,v^2,w^2$ sind.

Wenn man diese Discriminante aber nach der Formel (10), \S\,46 bildet, indem man
\[
 a_0=1,\quad a_1=2a,\quad a_2=a^2-4c,\quad a_3=-b^2
\]
setzt, so folgt
\begin{equation}
D=-4a^3b^2+144acb^2+16a^4c-128a^2c^2+256c^3-27b^4.
\tag{12}
\end{equation}
Wenn man dagegen zur Bildung der Discriminante der cubischen Gleichung (11) die Formel (8) anwendet, so erhält man
\begin{equation}
27D=4(a^2+12c)^3-(2a^3-72ac+27b^2)^2.
\tag{13}
\end{equation}
Man leitet hieraus die Discriminante der allgemeinen Function vierten Grades
\begin{equation}
f(x)=a_0x^4+a_1x^3+a_2x^2+a_3x+a_4
\tag{14}
\end{equation}
ab, indem man nach \S\,34, (4) $a,b,c$ durch
\[
 -\frac{3a_1^2}{8}+a_2,\qquad
 \frac{a_1^3}{8}-\frac{a_1a_2}{2}+a_3,
\qquad
 -\frac{3a_1^4}{256}+\frac{a_1^2a_2}{16}-\frac{a_1a_3}{4}+a_4
\]
ausdrückt, dann $a_1,a_2,a_3,a_4$ durch
\[
\frac{a_1}{a_0},\quad \frac{a_2}{a_0},\quad \frac{a_3}{a_0},\quad \frac{a_4}{a_0}
\]
ersetzt und mit $a_0^6$ multiplicirt. Setzt man nach Ausführung dieser einfachen Rechnung zur Abkürzung
\begin{align}
A&=a_2^2-3a_1a_3+12a_0a_4,
\tag{15}\\
B&=27a_1^2a_4+27a_0a_3^2+2a_2^3-72a_0a_2a_4-9a_1a_2a_3,
\tag{16}
\end{align}
so erhält man aus (13)
\begin{equation}
27D=4A^3-B^2.
\tag{17}
\end{equation}
Die Grössen $A$ und $B$, die uns im nächsten Abschnitt noch beschäftigen werden, heissen die erste und zweite Invariante der biquadratischen Function $f$. Man kann also nach (17) die Discriminante rational durch diese Grössen ausdrücken, jedoch nicht so, dass die Coefficienten ganzzahlig werden; führt man die in (17) vorgeschriebene Rechnung aus, um $D$ durch die Coefficienten $a_0,a_1,a_2,a_3,a_4$ selbst darzustellen, so muss sich der Factor $27$ noch fortheben lassen. Wir wollen den langen Ausdruck, dessen Berechnung keinerlei Schwierigkeit bietet, nicht hierher setzen.

\sect{48}{Resultanten}

Eine andere Anwendung der Theorie der symmetrischen Functionen, die übrigens die Discriminantenbildung als speciellen Fall enthält, ist die Bildung der Resultanten.

Es seien
\begin{align}
f(x)&=a_0x^n+a_1x^{n-1}+a_2x^{n-2}+\cdots+a_n,
\tag{1}\\
\varphi(x)&=b_0x^m+b_1x^{m-1}+b_2x^{m-2}+\cdots+b_m
\tag{2}
\end{align}
zwei ganze rationale Functionen vom $n$ten und $m$ten Grade und $\alpha_1,\alpha_2,\ldots,\alpha_n$ seien die Wurzeln der ersten, so dass auch
\begin{equation}
f(x)=a_0(x-\alpha_1)(x-\alpha_2)\cdots(x-\alpha_n)
\tag{3}
\end{equation}
gesetzt werden kann. Nun ist das Product
\[
\varphi(\alpha_1)\varphi(\alpha_2)\cdots\varphi(\alpha_n)
\]
eine symmetrische Function der Wurzeln von $f(x)$ und kann also ganz und rational durch
\[
\frac{a_1}{a_0},\quad \frac{a_2}{a_0},\ldots,\frac{a_n}{a_0}
\]
ausgedrückt werden. Die Ordnung dieser symmetrischen Function ist $(m,m\ldots m)$ (\S\,45) und daher ist
\begin{equation}
R=a_0^m\varphi(\alpha_1)\varphi(\alpha_2)\cdots\varphi(\alpha_n)
\tag{4}
\end{equation}
eine ganze rationale und homogene Function $m$ter Ordnung von $a_0,a_1,\\ldots,a_n$. Sie ist auch, wie ihr Ausdruck zeigt, eine ganze homogene Function $n$ter Ordnung von $b_0,b_1,\ldots,b_m$. Sie wird die Resultante der beiden Functionen $f(x)$ und $\varphi(x)$ genannt. Diese Bezeichnung rechtfertigt sich durch eine zweite Darstellung.

Bezeichnen wir die Wurzeln von $\varphi(x)$ mit $\beta_1,\beta_2,\ldots,\beta_m$, setzen also
\begin{equation}
\varphi(x)=b_0(x-\beta_1)(x-\beta_2)\cdots(x-\beta_m),
\tag{5}
\end{equation}
so wird nach (4)
\begin{equation}
R=a_0^m b_0^n \prod_{i,k}(\alpha_i-\beta_k),
\tag{6}
\end{equation}
worin das Productzeichen $\prod$ sich auf alle Indices $i=1,2,\ldots,n$, $k=1,2,\ldots,m$ bezieht.

Dieser Ausdruck zeigt nun, dass $R$, abgesehen von dem Factor $(-1)^{mn}$, in der gleichen Weise von $\varphi(x)$ abhängt, wie von $f(x)$, und dass wir also auch setzen können
\begin{equation}
(-1)^{mn}R=b_0^n f(\beta_1)f(\beta_2)\cdots f(\beta_m).
\tag{7}
\end{equation}
Die Discriminanten sind hiernach ein specieller Fall der Resultanten, den man erhält, wenn man für $\varphi(x)$ die Derivirte $f'(x)$ nimmt.

Das Verschwinden der Resultante von $f(x)$ und $\varphi(x)$ ist die nothwendige und hinreichende Bedingung dafür, dass die beiden Gleichungen
\[
f(x)=0\quad\text{und}\quad \varphi(x)=0
\]
eine gemeinsame Wurzel haben.

Die Gleichung, die durch Nullsetzen der Resultante entsteht, wird auch die Endgleichung aus $f(x)=0$ und $\varphi(x)=0$ genannt und die Aufstellung der Endgleichung die Elimination von $x$.

Zur Berechnung der Resultanten kann man den Algorithmus des grössten gemeinsamen Theilers anwenden (\S\,6). Ist nämlich $m\le n$, so kann man eine ganze Function $Q$ und eine ebensolche Function $\psi(x)$, deren Grad kleiner ist als $m$, so bestimmen, dass
\[
f(x)=Q\varphi(x)+\psi(x),
\]
also, wenn man für $x$ eine der Wurzeln $\beta_i$ setzt,
\[
f(\beta_i)=\psi(\beta_i),
\]
also
\begin{equation}
b_0^n\prod_i f(\beta_i)=b_0^n\prod_i\psi(\beta_i);
\tag{8}
\end{equation}
ist $m'$ der Grad von $\psi(x)$, so ist
\begin{equation}
b_0^{m'}\prod_i\psi(\beta_i)=R'
\tag{9}
\end{equation}
die Resultante von $\varphi(x)$ und $\psi(x)$ und daher
\begin{equation}
R=b_0^{n-m'}R'.
\tag{10}
\end{equation}
Die Bildung der Resultante $R$ ist hierdurch auf die Bildung von $R'$ zurückgeführt. Man kann nun $\varphi(x)$ wieder durch $\psi(x)$ theilen, und führt so die Bildung von $R'$ auf die Bildung einer Resultante von Functionen immer niedrigeren Grades zurück, bis man schliesslich auf eine Constante, d. h. eine Function der Coefficienten $a,b$ kommt, die mit der gesuchten Resultante identisch sein muss.

Die Berechnung der Resultanten ist meist sehr weitläufig. Wir wollen nur das einfachste Beispiel der Resultante zweier quadratischer Functionen anführen:
\[
f(x)=a_0x^2+a_1x+a_2,
\qquad
\varphi(x)=b_0x^2+b_1x+b_2.
\]
Man kann hier direct das Product
\[
R=a_0^2(b_0\alpha_1^2+b_1\alpha_1+b_2)(b_0\alpha_2^2+b_1\alpha_2+b_2)
\]
berechnen, und erhält, wenn man
\[
-a_0(\alpha_1+\alpha_2)=a_1,\qquad a_0\alpha_1\alpha_2=a_2
\]
setzt,
\begin{equation}
R=a_0^2b_2^2+a_2^2b_0^2-b_0b_1a_1a_2-a_0a_1b_1b_2+b_0b_2a_1^2+a_0a_2b_1^2-2a_0b_0a_2b_2,
\tag{11}
\end{equation}
oder was damit identisch ist
\begin{equation}
R=(a_0b_2-a_2b_0)^2-(a_0b_1-a_1b_0)(a_1b_2-a_2b_1).
\tag{12}
\end{equation}
Die einzelnen Glieder der Resultante der zwei Functionen $f(x)$ und $\varphi(x)$ vom $n$ten und $m$ten Grade haben, von einem numerischen (ganzzahligen) Factor abgesehen, die Gestalt
\begin{equation}
a_0^{\nu_0}a_1^{\nu_1}\cdots a_n^{\nu_n} b_0^{\mu_0}b_1^{\mu_1}\cdots b_m^{\mu_m},
\tag{13}
\end{equation}
worin, wie aus den oben bestimmten Graden hervorgeht,
\begin{equation}
\nu_0+\nu_1+\cdots+\nu_n=m,\qquad
\mu_0+\mu_1+\cdots+\mu_m=n.
\tag{14}
\end{equation}
Wir können aber noch eine andere Relation zwischen diesen Exponenten angeben. Da nämlich [nach (6)] $R$ eine homogene Function $nm$ten Grades von den $n+m$ Variablen $\alpha,\beta$ ist, da ferner $a_0$ vom nullten, $a_1$ vom ersten, $a_2$ vom zweiten etc. Grade in den $\alpha$, allgemein $a_k$ und $b_k$ vom $k$ten Grade in diesen Variablen sind, so folgt
\begin{equation}
\nu_1+2\nu_2+3\nu_3+\cdots+n\nu_n+
\mu_1+2\mu_2+3\mu_3+\cdots+m\mu_m=nm,
\tag{15}
\end{equation}
eine Relation, aus der sich ein wichtiger Schluss ziehen lässt, den wir im folgenden Paragraphen etwas ausführlicher besprechen wollen.

\sect{49}{Elimination. Theorem von Bezout}

Wenn in den beiden Functionen
\begin{equation}
\begin{aligned}
f(x)&=a_0x^n+a_1x^{n-1}+\cdots+a_n,\\
\varphi(x)&=b_0x^m+b_1x^{m-1}+\cdots+b_m
\end{aligned}
\tag{1}
\end{equation}
die Coefficienten $a,b$ selbst wieder ganze rationale Functionen einer Grösse $y$ sind, so wird auch die Resultante eine ganze rationale Function von $y$ werden, die wir mit $F(y)$ bezeichnen wollen. Die Bildung dieser Gleichung ist die Elimination von $x$ aus den beiden Gleichungen (1).

Die Gleichung
\begin{equation}
F(y)=0
\tag{2}
\end{equation}
hat alle die Werthe von $y$ zu Wurzeln, wofür die beiden Gleichungen
\begin{equation}
f(x)=0,\qquad \varphi(x)=0
\tag{3}
\end{equation}
eine gemeinsame Wurzel haben. Ist die Bedingung (2) befriedigt, so findet man die Werthe von $x$, die den zwei Gleichungen (3) genügen, indem man den grössten gemeinschaftlichen Theiler von $f(x)$ und $\varphi(x)$ aufsucht. Dieser gemeinschaftliche Theiler ist vom ersten Grade und bestimmt also $x$ rational, wenn nur eine solche gemeinschaftliche Wurzel vorhanden ist; im anderen Falle wird er von entsprechend höherem Grade, und die Bestimmung von $x$ erfordert noch die Lösung einer Gleichung höheren Grades.

Gehört zu jeder Wurzel der Resultante nur eine gemeinschaftliche Wurzel $x$, so ist die Zahl der Werthepaare $x,y$, die den Gleichungen (3) genügen, gleich dem Grade der Resultante in Bezug auf $y$. Sind also die Coefficienten $a$ vom Grade $\nu$, $b$ vom Grade $\mu$ in Bezug auf $y$, so ist nach der Gradbestimmung des vorigen Paragraphen der Grad der Resultante
\begin{equation}
n\nu+m\mu;
\tag{5}
\end{equation}
so gross ist also die Anzahl der gemeinsamen Wurzelpaare. Diese Zahl kann aber noch dadurch modificirt werden, dass möglicherweise die Coefficienten so beschaffen sein können, dass die höchste Potenz von $y$ aus der Endgleichung wegfällt. Man würde dann die Uebereinstimmung der Zahl mit (5) nur durch die willkürliche Hinzufügung unendlich vieler Wurzeln retten können. Auch mehrfache Wurzeln geben zu Bedenken Anlass. Man begegnet diesen Uebelständen theilweise dadurch, dass man die homogene Form der Gleichungen zu Grunde legt, eine Form, in der sich das Eliminationsproblem besonders in der Geometrie einstellt.

Wenn nämlich die Function $f(x)$ aus einer homogenen Function $n$ter Ordnung der drei Variablen $x,y,z$ (einer ternären Form) durch Ordnen nach Potenzen von $x$ hervorgegangen ist, so sind die Coefficienten $a_0,a_1,a_2,\ldots,a_n$ homogene Formen der beiden Variablen $y,z$, von dem Grade, den der Index angiebt, also $a_0$ eine Constante, $a_1$ eine lineare, $a_2$ eine quadratische Form u. s. f., und Entsprechendes gilt von den Coefficienten $b_0,b_1,\ldots,b_m$ der Function $\varphi(x)$.

Bedeuten $x,y,z$ Dreieckscoordinaten in der Ebene, so sind $f(x)=0$, $\varphi(x)=0$ die Gleichungen einer Curve $n$ter und $m$ter Ordnung, und die gemeinsamen Wurzeln dieser beiden Gleichungen (die Verhältnisse $x:y:z$) bedeuten die Schnittpunkte der beiden Curven. Die Endgleichung, die man durch Elimination von $x$ erhält, ist eine homogene Gleichung, deren Grad sich aus (15), \S\,48 gleich $mn$ ergiebt.

Eine Verminderung des Grades kann hier nicht stattfinden; immer ist die Endgleichung eine homogene Gleichung des $nm$ten Grades für die beiden Unbekannten $y,z$. Die einzige Ausnahme, auf die wir zu achten ist, ist die, dass die Resultante identisch (für alle $y,z$) verschwindet, dass also unendlich viele gemeinsame Wurzeln vorhanden sind. Dieser Fall hat in der Geometrie die Bedeutung, dass die beiden Curven einen Curventheil und nicht bloss einzelne Punkte gemein haben.

Wenn die Gleichungen $f=0$, $\varphi=0$ nur eine endliche Anzahl gemeinsamer Wurzeln haben, so können wir in den Ausdrücken
\[
y_1=y+\beta x,
\qquad
z_1=z+\gamma x
\]
die Constanten $\beta,\gamma$ so einrichten, dass, wenn einem Werth des Verhältnisses $y:z$ mehrere verschiedene der Gleichungen (1) genügende Werthe von $x$ entsprechen, die zugehörigen Werthe von $y_1:z_1$ von einander verschieden ausfallen.

Wir können dann $f$ und $\varphi$ auch als homogene Functionen vom Grade $m$ und $n$ der Veränderlichen $x,y_1,z_1$ ansehen, und die Resultante wird dann eine homogene Function $mn$ten Grades $R(y_1,z_1)$ von $y_1$ und $z_1$. Setzen wir darin
\[
y_1=b'+b\lambda,
\qquad
z_1=c'+c\lambda,
\]
so erhalten wir eine Gleichung $mn$ten Grades in $\lambda$, und wir können, wenn $R(y_1,z_1)$ nicht identisch verschwindet, die Constanten $b,b',c,c'$ so annehmen, dass der Coefficient der höchsten Potenz von $\lambda$, nämlich $R(b,c)$, von Null verschieden ist. Dann entspricht jeder Wurzel der Resultante nur ein Werthsystem des Verhältnisses $x:y:z$, und die Anzahl der gemeinschaftlichen Wurzeln von $f$ und $\varphi$ ist also höchstens gleich $mn$. Um den Satz aussprechen zu können, dass die Zahl der gemeinsamen Wurzeln immer gleich $mn$ ist, muss man noch ein Uebereinkommen treffen, wonach gewisse Werthepaare mehrfach zu zählen sind.

Der Satz, der, in geometrischem Gewande, besagt, dass sich zwei Curven $m$ter und $n$ter Ordnung in $mn$ Punkten schneiden, wird das Bezout'sche Theorem genannt.

\sect{50}{Elimination aus mehreren Gleichungen}

Die Principien der Elimination, die wir im Vorhergehenden auf zwei Gleichungen angewandt haben, lassen sich ohne Schwierigkeit auf mehrere Gleichungen mit einer entsprechenden Anzahl von Unbekannten ausdehnen.

Nehmen wir zunächst die Gleichungen
\begin{equation}
f(x)=0,\qquad \varphi(x)=0,\qquad \psi(x)=0,
\tag{1}
\end{equation}
der Grade $m,n,p$, deren Coefficienten $a_i,b_i,c_i$ zunächst als unabhängige Veränderliche betrachtet werden sollen.

Wir wenden auf zwei von ihnen, etwa auf $f(x),\varphi(x)$, den Algorithmus des grössten gemeinschaftlichen Theilers an und erhalten einen letzten Rest, der eine rationale Function der $a$ und $b$ ist und der, von Nennern befreit, mit der Resultante $R$ von $f(x)$ und $\varphi(x)$ übereinstimmt, und einen vorletzten, der eine lineare Function von $x$ ist. Die Bedingung für das Vorhandensein einer gemeinsamen Wurzel von $f=0$ und $\varphi=0$ ist das Verschwinden der Resultante, und der vorletzte Rest giebt unter dieser Voraussetzung für die gemeinsame Wurzel einen rationalen Ausdruck durch die Coefficienten.

Setzt man diesen Ausdruck in $\psi(x)$ ein, so erhält man, wenn man alle Nenner fortschafft, eine ganze rationale Function der Coefficienten $a,b,c$, die wir mit $S$ bezeichnen wollen, und deren Verschwinden in Verbindung mit $R=0$ die nothwendige und hinreichende Bedingung dafür enthält, dass die drei Gleichungen (1) eine gemeinsame Wurzel haben.

Ersetzen wir nun die Coefficienten $a,b,c$ der Functionen (1) durch ganze rationale Functionen einer zweiten Unbekannten $y$, deren Coefficienten mit $\alpha,\beta,\gamma$ bezeichnet und vorläufig als unabhängige Veränderliche betrachtet werden mögen. $R$ und $S$ werden dann ganze rationale Functionen von $\alpha,\beta,\gamma$ und $y$, und wenn wir also die Gleichungen
\begin{equation}
R=0,\qquad S=0
\tag{2}
\end{equation}
als Gleichungen für $y$ betrachten, so können wir nach den Regeln des vorigen Paragraphen ihre Resultante $T$ bilden, die, wenn sie nicht identisch verschwindet, eine ganze rationale Function der Coefficienten $\alpha,\beta,\gamma$ sein wird, und die Gleichung
\begin{equation}
T=0
\tag{3}
\end{equation}
ist die nothwendige und hinreichende Bedingung dafür, dass das Gleichungssystem (1) durch ein oder mehrere Werthepaare $x,y$ befriedigt werden kann. $T$ heisst die Resultante des Systems (1). Man kann also aus den drei Gleichungen (1) durch Elimination der Unbekannten $x$ und $y$ eine Endgleichung für $z$ bilden.

Ersetzt man nun die $\alpha,\beta,\gamma$ wieder durch ganze rationale Functionen einer dritten Unbekannten $z$, so geht $T$ in eine Gleichung für $z$ über. Die Functionen $f,\varphi,\psi$ werden Functionen der drei Unbekannten $x,y,z$. Der Grad der Gleichung (3) giebt die Anzahl der Werthsysteme $x,y,z$, für die die drei Gleichungen (1) zugleich befriedigt sind.

Sind aber die Coefficienten so beschaffen, dass $T$ identisch (für alle Werthe von $z$) verschwindet, so giebt es unendlich viele Werthsysteme $x,y,z$, die den Gleichungen (1) zugleich genügen. Betrachtet man $x,y,z$ als Cartesische Coordinaten, so stellen die Gleichungen (1), jede für sich, eine Oberfläche dar. Der Grad der Endgleichung $T=0$ in Bezug auf $z$ giebt die Anzahl der Schnittpunkte dieser drei Oberflächen an. Ist aber $T$ identisch Null, so haben die drei Flächen entweder einen Flächentheil oder eine Curve gemeinschaftlich.

Den Grad der Endgleichung können wir auf folgendem Wege bestimmen, der gar keine weitläufigen Betrachtungen über die Grade der Functionen $R,S,T$ erfordert. Wir führen, um uns von Ausnahmefällen unabhängig zu machen, homogene Variable ein, indem wir eine vierte, $u$, hinzufügen, und also drei homogene quaternäre Gleichungen
\begin{equation}
F(x,y,z,u)=0,
\qquad
\Phi(x,y,z,u)=0,
\qquad
\Psi(x,y,z,u)=0
\tag{4}
\end{equation}
von den Ordnungen $m,n,p$ betrachten.

Nehmen wir den besonderen Fall an, dass jede dieser drei Functionen in lauter lineare Factoren zerfällt, so ist klar, dass, wenn die Anzahl der gemeinsamen Lösungen nicht unendlich ist, ihre Anzahl gleich dem Product $mnp$ sein muss; denn wir können je einen linearen Factor von $F,\Phi,\Psi$ gleichzeitig Null setzen und erhalten so $mnp$ Systeme von je drei linearen Gleichungen, die, wenn sie von einander unabhängig sind, ebenso viele Werthsysteme für die Verhältnisse $x:y:z:u$ ergeben.

Der Schluss, dass im allgemeinen Fall der Grad der Endgleichung von (4) ebenso gross ist, ist auf den ersten Blick bedenklich. Das Bedenken kann aber nach einer Bemerkung von F. Klein durch folgende einfache Ueberlegung vollständig beseitigt werden.

Denken wir uns die Endgleichung aus den Gleichungen (4) zunächst unter der Voraussetzung gebildet, dass die Coefficienten in den Gleichungen (4) von einander unabhängige Variable sind, so wird die Resultante, die sich nach Elimination von $x,y$ ergiebt, gewiss nicht identisch Null, da sie schon für specielle Coefficientensysteme, z. B. wenn $F,\Phi,\Psi$ in lineare Factoren zerfallen, nicht identisch verschwindet, und wir erhalten eine homogene Endgleichung für die zwei Unbekannten $z,u$ von einem gewissen Grade, der zu ermitteln ist.

Specialisirt man die Coefficienten irgend wie, so kann zwar möglicherweise die Endgleichung identisch befriedigt werden; wenn aber dieser Fall nicht eintritt, so kann ihre Ordnung nicht geändert werden (da ja alle Glieder von derselben Ordnung sind). Da nun durch die oben angenommene Specialisirung, die sich in dem Zerfallen der Functionen $F,\Phi,\Psi$ ausspricht, das identische Verschwinden nicht stattfindet, sondern eine Endgleichung von der Ordnung $mnp$ auftritt, so muss dies auch der Grad der Endgleichung im allgemeinen Falle sein.

Hiermit ist das Bezout'sche Theorem für den Fall von drei homogenen quaternären Gleichungen bewiesen.

Dass wir uns hier auf drei Gleichungen beschränkt haben, ist lediglich im Interesse der Einfachheit geschehen. Die Erweiterung, die für den Fall von $n$ Gleichungen mit $n+1$ homogenen Unbekannten erforderlich ist, ergiebt sich von selbst, und kann dem Leser überlassen bleiben.

Für unsere allgemeine Aufgabe wollen wir nur noch bemerken, dass sich aus diesen Betrachtungen ergiebt, dass das Problem der Algebra, sofern es sich auf die Lösung irgend welcher algebraischer Gleichungen bezieht, damit zurückgeführt ist auf die Behandlung einer Kette von Gleichungen mit je einer Unbekannten.

\sect{51}{Zerlegbare und unzerlegbare Functionen}

Unter den ganzen rationalen Functionen mehrerer Veränderlichen müssen zerlegbare und unzerlegbare unterschieden werden.

Eine ganze rationale Function von irgend welchen Veränderlichen heisst zerlegbar, wenn sie als Product von wenigstens zwei ganzen rationalen Functionen derselben Veränderlichen darstellbar ist; ist sie nicht als ein solches Product darstellbar, so heisst sie unzerlegbar.

Wenn eine Function zerlegbar ist, so ist der Grad jedes der Factoren niedriger als der Grad der Function selbst. Eine lineare Function ist also immer unzerlegbar, und jede ganze rationale Function lässt sich in eine endliche Anzahl unzerlegbarer Functionen zerlegen.

Wir nennen eine ganze Function $W$ durch eine andere $w$ theilbar, wenn eine dritte ganze Function $w'$ existirt, so dass
\[
W=ww'
\]
ist. Daraus folgt dann, dass, wenn $U$ eine durch $w$ theilbare und $V$ eine beliebige ganze Function ist, auch das Product $UV$ durch $w$ theilbar ist, und dass, wenn $U,U',U'',\ldots$ durch $w$ theilbare, $V,V',V'',\ldots$ beliebige ganze rationale Functionen sind, auch $UV+U'V'+U''V''+\cdots$ durch $w$ theilbar ist.

Ist $W$ durch $w$ theilbar, so sagt man auch, $w$ geht in $W$ auf.

Zwei ganze rationale Functionen $U,V$, die nicht durch eine und dieselbe ganze rationale Function theilbar sind, heissen relativ prim oder theilerfremd.

Wir beweisen den Satz:

\begin{enumerate}[label=\Roman*.,leftmargin=2.4em]
\item Sind $U,V,v$ ganze rationale Functionen irgend welcher Veränderlichen, sind $U$ und $v$ relativ prim und $UV$ durch $v$ theilbar, so ist $V$ durch $v$ theilbar.
\end{enumerate}

Dieser Satz entspricht genau einem bekannten Fundamentalsatz aus der Lehre von den ganzen Zahlen, dass nämlich, wenn ein Product von zwei ganzen Zahlen durch eine dritte ganze Zahl theilbar ist, die zu dem einen Factor theilerfremd ist, der andere Factor durch diese Zahl theilbar sein muss.

Wir beweisen ihn durch vollständige Induction. Sind $U,V,v$ nur von einer Veränderlichen $x$ abhängig, so ist der Satz richtig; denn nach \S\,6 (Satz I) kann man in diesem Falle, wenn $U$ und $v$ relativ prim sind, zwei andere ganze Functionen $P$ und $p$ von $x$ so bestimmen, dass
\[
PU+pv=1,
\]
woraus durch Multiplication mit $V$
\[
PUV+pVv=V
\]
folgt, und daraus ersieht man, dass, wenn $UV$ durch $v$ theilbar ist, auch $V$ durch $v$ theilbar sein muss.

Wir nehmen also an, der Satz, den wir beweisen wollen, sei für Functionen von $n$ und weniger Veränderlichen bewiesen, und wir leiten daraus seine Richtigkeit für Functionen von $n+1$ Veränderlichen ab.

Dazu ist erforderlich, dass wir aus dem als richtig vorausgesetzten Theorem I einige Folgerungen ziehen, als deren Schluss sich dann die Gültigkeit des Theorems für die nächst höhere Variablenzahl ergiebt.

Ist $v$ eine unzerlegbare Function, so ist eine andere Function $U$ derselben Veränderlichen entweder relativ prim zu $v$ oder durch $v$ theilbar. Daraus folgt nach I., dass ein Product von zwei Functionen $UV$ nur dann durch $v$ theilbar sein kann, wenn einer der beiden Factoren durch $v$ theilbar ist. Dasselbe gilt für ein Product von mehreren Functionen, und so folgt aus I. das Theorem:

\begin{enumerate}[label=\Roman*.,leftmargin=2.4em,start=2]
\item Ein Product aus mehreren ganzen rationalen Functionen ist nur dann durch eine unzerlegbare Function $v$ theilbar, wenn wenigstens einer der Factoren des Productes durch $v$ theilbar ist.
\end{enumerate}

Wenn eine ganze rationale Function $U$ auf zwei Arten in unzerlegbare Factoren zerlegt ist,
\[
U=vv'v''\cdots=ww'w''\cdots,
\]
so muss nach II. wenigstens einer der Factoren $v,v',v'',\ldots$ durch $w$ theilbar sein, also etwa $v$. Dann aber kann, da auch $v$ unzerlegbar ist, $v$ von $w$ nur durch einen constanten Factor verschieden sein.

Demnach ist, wenn $c$ dieser constante Factor ist,
\[
cv'v''\cdots=w'w''\cdots,
\]
woraus folgt, dass eine der Functionen $v',v'',\ldots$, etwa $v'$, durch $w'$ theilbar ist, und sich also von $w'$ nur durch einen constanten Factor unterscheidet.

Wir erhalten also als zweite Folgerung aus dem Theorem I:

\begin{enumerate}[label=\Roman*.,leftmargin=2.4em,start=3]
\item Eine ganze rationale Function kann, von constanten Factoren abgesehen, nur auf eine Art in unzerlegbare Factoren zerlegt werden.
\end{enumerate}

Hieraus ergiebt sich der Begriff des grössten gemeinschaftlichen Theilers von zwei oder mehr ganzen rationalen Functionen $U,V,\ldots$. Man versteht darunter das Product aller unzerlegbaren Factoren, die in den Zerlegungen jeder der Functionen $U,V,\ldots$ vorkommen, oder die Function möglichst hohen Grades, die in allen Functionen $U,V,\ldots$ aufgeht. Nach III. ist diese Function, von einem constanten Factor abgesehen, für jedes Functionensystem $U,V,\ldots$ vollständig bestimmt.

Mehrere Functionen $U,V,W,\ldots$ heissen relativ prim, wenn keine zwei von ihnen einen gemeinsamen Theiler haben.

Alle diese Definitionen und Sätze sind genau analog mit sehr bekannten Sätzen der elementaren Zahlenlehre, die dort als Folgerungen des Algorithmus des grössten gemeinschaftlichen Theilers auftreten, nur dass hier die unzerlegbaren Functionen die Rolle der Primzahlen übernehmen.

Wir führen noch einen solchen Satz an:

\begin{enumerate}[label=\Roman*.,leftmargin=2.4em,start=4]
\item Sind $u,v$ relativ prim und $Uu,Uv$ durch $w$ theilbar, so ist $U$ durch $w$ theilbar.
\end{enumerate}

Denn zerlegt man die ganzen rationalen Functionen $U,u,v,w$ in ihre unzerlegbaren Factoren, so muss irgend ein Factor von $w$, da er nicht in $u$ und $v$ zugleich vorkommen kann, in $U$ aufgehen. Hebt man ihn aus $U$ und $w$ weg, so kann man ebenso für einen nächsten Factor von $w$ schliessen u. s. f.

Wir betrachten nun ganze rationale Functionen einer Veränderlichen $t$
\[
f(t)=u_0t^m+u_1t^{m-1}+\cdots+u_{m-1}t+u_m,
\]
deren Coefficienten ganze rationale Functionen von $n$ Veränderlichen $x$ sind, von denen $t$ unabhängig ist.

Sind die Coefficienten $u_0,u_1,\ldots,u_m$ ohne gemeinsamen Theiler, so heisst $f(t)$ primitiv, im anderen Falle wird der grösste gemeinschaftliche Theiler der Coefficienten $u_0,u_1,\ldots,u_m$ der Theiler der Function $f(t)$ genannt (vgl. \S\,2).

Wir schliessen, immer unter Voraussetzung der Gültigkeit von I:

\begin{enumerate}[label=\Roman*.,leftmargin=2.4em,start=5]
\item Das Product von zwei primitiven Functionen $f(t)$ und $\varphi(t)$ ist wieder eine primitive Function und der Theiler eines Productes zweier imprimitiver Functionen ist gleich dem Product der Theiler beider Factoren.
\end{enumerate}

Der Beweis ist ganz so wie für den entsprechenden Satz in \S\,2.

Es seien
\begin{equation}
\begin{aligned}
f(t)&=u_0t^m+u_1t^{m-1}+\cdots+u_m,\\
\varphi(t)&=v_0t^\mu+v_1t^{\mu-1}+\cdots+v_\mu
\end{aligned}
\tag{1}
\end{equation}
zwei primitive Functionen und
\begin{equation}
F(t)=U_0t^{m+\mu}+U_1t^{m+
\mu-1}+\cdots+U_{m+\mu}
\tag{2}
\end{equation}
ihr Product. Es ist dann, wenn $r$ und $s$ irgend zwei Ziffern aus den Reihen $0,1,\ldots,m$ und $0,1,2,\ldots,\mu$ bedeuten (\S\,2),
\begin{equation}
U_{r+s}=u_rv_s+u_{r-1}v_{s+1}+\cdots+u_{r+1}v_{s-1}+\cdots.
\tag{3}
\end{equation}
Wenn nun $w$ irgend eine unzerlegbare Function ist, die weder in allen $u_r$ noch in allen $v_s$ aufgeht, so wählen wir in (3) $r$ und $s$ so, dass $u_r$ das erste nicht durch $w$ theilbare $u$ ist, also $u_{r-1},u_{r-2},\ldots$ durch $w$ theilbar sind, und dass ebenso $v_s$ das erste, nicht durch $w$ theilbare $v$ wird. Dann kann auch $U_{r+s}$ nicht durch $w$ theilbar sein, weil alle Glieder, mit Ausnahme des ersten, durch $w$ theilbar sind, d. h. $F(t)$ ist primitiv.

Daraus folgt unmittelbar, wenn $p$ und $q$ irgend welche ganze rationale Functionen der $x$ sind, dass $pq$ der Theiler des Productes der beiden imprimitiven Functionen $pf(t),q\varphi(t)$ ist, also der zweite Theil des Satzes V. Daraus folgt weiter:

\begin{enumerate}[label=\Roman*.,leftmargin=2.4em,start=6]
\item Wenn eine ganze rationale Function $F(t)$ der $n+1$ Variablen $x$ und $t$ in zwei Factoren zerlegbar ist, die in Bezug auf $t$ ganz, in Bezug auf die $x$ wenigstens rational sind, so ist sie auch in zwei Functionen zerlegbar, die in $x$ und $t$ ganz und rational sind.
\end{enumerate}

Denn nach der Voraussetzung giebt es eine ganze rationale Function $w$ der $x$ allein und zwei ganze rationale Functionen der $x$ und $t$, $f_1(t),\varphi_1(t)$, so dass
\[
wF(t)=f_1(t)\varphi_1(t),
\]
und nach (V) muss $w$ in dem Product der Theiler von $f_1(t)$ und $\varphi_1(t)$ aufgehen, so dass wir auch
\[
F(t)=f(t)\varphi(t)
\]
erhalten, worin $f(t),\varphi(t)$ gleichfalls ganze rationale Functionen von $x$ und $t$ und zwar in Bezug auf $t$ von demselben Grade wie $f_1(t)$ und $\varphi_1(t)$ sind, w. z. b. w.

Hieraus schliessen wir weiter, dass zwei Functionen $F(t),f(t)$, die als ganze rationale Functionen der $n+1$ Veränderlichen $x$ und $t$ betrachtet, in dem oben definirten Sinne relativ prim sind, sich auch, als Functionen von $t$ allein betrachtet, und nach dem Algorithmus des grössten gemeinschaftlichen Theilers behandelt (\S\,6), als relativ prim erweisen müssen. Denn wenn sie einen gemeinsamen Theiler hätten, der in Bezug auf $t$ ganz, in Bezug auf die $x$ gebrochen wäre, so liesse sich eine ganze Function $T$ von $x$ und $t$ und eine ganze Function $P$ der $x$ allein so bestimmen, dass
\[
PF(t)=TF_1(t),\qquad Pf(t)=Tf_1(t)
\]
wäre, worin $F_1,f_1$ ganze Functionen ohne gemeinsamen Theiler sind. Nach IV müsste also $P$ im Theiler von $T$ aufgehen, und es könnten $F(t)$ und $f(t)$ nicht relativ prim sein. Es lassen sich also nach \S\,6 zwei ganze rationale Functionen $Q,q$ von $x$ und $t$ und eine ganze rationale Function $X$ von den $x$ allein so bestimmen, dass die Identität
\begin{equation}
QF(t)+qf(t)=X
\tag{4}
\end{equation}
besteht.

Ist nun $\Phi(t)$ eine weitere ganze rationale Function von $x$ und $t$, so folgt aus (4) durch Multiplication mit $\Phi(t)$
\[
Q\Phi(t)F(t)+q\Phi(t)f(t)=X\Phi(t),
\]
und wenn also $\Phi(t)F(t)$ durch $f(t)$ theilbar ist, so ist hiernach auch $X\Phi(t)$ durch $f(t)$ theilbar.

Demnach können wir, wenn $\varphi(t)$ und $\psi(t)$ wieder zwei ganze Functionen von $x$ und $t$ bedeuten, setzen
\begin{equation}
\begin{aligned}
X\Phi(t)&=\varphi(t)f(t),\\
F(t)\Phi(t)&=\psi(t)f(t).
\end{aligned}
\tag{5}
\end{equation}
Multiplicirt man die zweite dieser Gleichungen mit $X$ und setzt aus der ersten für $X\Phi(t)$ den Ausdruck $\varphi(t)f(t)$ ein, so lässt sich $f(t)$ wegheben und es folgt
\[
\varphi(t)F(t)=X\psi(t).
\]
Nach IV. muss also $X$ sowohl im Theiler von $\varphi(t)f(t)$ als im Theiler von $\varphi(t)F(t)$ aufgehen, und da die Functionen $f(t)$ und $F(t)$ und mithin auch ihre Theiler relativ prim sind, so muss $X$ im Theiler von $\varphi(t)$ aufgehen. Setzen wir also demnach
\[
\varphi(t)=X\varphi_1(t),
\]
so folgt aus (5)
\[
\Phi(t)=\varphi_1(t)f(t),
\]
d. h. $\Phi(t)$ ist durch $f(t)$ theilbar, also:

\begin{enumerate}[label=\Roman*.,leftmargin=2.4em,start=7]
\item Sind $F(t)$ und $f(t)$ relativ prim und $\Phi(t)F(t)$ durch $f(t)$ theilbar, so ist $\Phi(t)$ durch $f(t)$ theilbar.
\end{enumerate}

Dies aber ist nichts Anderes als das Theorem I für Functionen von $n+1$ Variablen, und I. ist somit allgemein bewiesen.

\sect{52}{Tschirnhausen-Transformation}

Wir machen von der Theorie der symmetrischen Functionen und der Resultantenbildung noch eine wichtige Anwendung auf die von Tschirnhausen\footnote{Acta eruditorum, Leipzig 1683.} herrührende Transformation algebraischer Gleichungen.

Es sei
\begin{equation}
f(x)=x^n+a_1x^{n-1}+a_2x^{n-2}+\cdots+a_{n-1}x+a_n
\tag{1}
\end{equation}
irgend eine Function $n$ten Grades, und
\begin{equation}
y=\frac{\chi(x)}{\psi(x)}
\tag{2}
\end{equation}
irgend eine (ganze oder gebrochene) rationale Function von $x$, bei der wir nur die eine Voraussetzung machen, dass $f(x)$ und $\psi(x)$ keinen gemeinsamen Theiler haben sollen. Nach dem Schlusssatz von \S\,6 können wir dann die Functionen $F(x)$ und $\varphi(x)$, letztere höchstens vom Grade $n-1$, so bestimmen, dass
\begin{equation}
\chi(x)=F(x)f(x)+\varphi(x)\psi(x)
\tag{3}
\end{equation}
wird, so dass, sobald für $x$ eine Wurzel der Gleichung
\begin{equation}
f(x)=0
\tag{4}
\end{equation}
gesetzt wird,
\[
\frac{\chi(x)}{\psi(x)}=\varphi(x)
\]
wird. Da wir nun die Werthe von $y$ nur für solche Werthe von $x$, die Wurzeln von (4) sind, betrachten werden, so verlieren wir nichts an Allgemeinheit, wenn wir von vornherein
\[
y=\varphi(x),
\]
also gleich einer ganzen Function $(n-1)$ten Grades setzen. Es sei also
\begin{equation}
y=\varphi(x)=\alpha_0+\alpha_1x+\alpha_2x^2+\cdots+\alpha_{n-1}x^{n-1},
\tag{5}
\end{equation}
worin die Coefficienten $\alpha$ vorläufig noch ganz unbestimmt bleiben.

Bezeichnen wir mit
\begin{equation}
x_1,x_2,\ldots,x_n
\tag{6}
\end{equation}
die Wurzeln der Gleichung (4), so ergeben sich aus (5) die entsprechenden Werthe von $y$:
\begin{equation}
y_1=\varphi(x_1),\qquad y_2=\varphi(x_2),\ldots, y_n=\varphi(x_n),
\tag{7}
\end{equation}
und das Product
\begin{equation}
\Phi(y)=(y-y_1)(y-y_2)\cdots(y-y_n)
\tag{8}
\end{equation}
ist eine symmetrische Function der Grössen (6), die sich also rational durch die Coefficienten von $f(x)$ ausdrücken lässt.

In dieser Weise dargestellt, möge $\Phi(y)$ den Ausdruck haben
\begin{equation}
\Phi(y)=y^n+b_1y^{n-1}+b_2y^{n-2}+\cdots+b_{n-1}y+b_n.
\tag{9}
\end{equation}
Die Coefficienten $b_\nu$ werden nach \S\,7 durch die $y_i$ bestimmt mittelst der Formeln
\[
b_1=-\sum y_1,
\qquad
b_2=\sum y_1y_2,
\qquad
b_3=-\sum y_1y_2y_3,\ldots,
\]
woraus sich ergiebt, dass $b_\nu$ eine ganze rationale und homogene Function $\nu$ten Grades von den Variablen $\alpha_0,\alpha_1,\ldots,\alpha_{n-1}$ ist.

$\Phi(y)$ ist nichts Anderes als die Resultante der beiden Gleichungen
\begin{equation}
f(x)=0,\qquad \varphi(x)-y=0,
\tag{10}
\end{equation}
und nach \S\,49 kann man auch umgekehrt die gemeinsame Wurzel dieser beiden Gleichungen (10), wenn $y$ einem der Werthe $y_1,y_2,\ldots,y_n$ gleich wird, rational durch die Coefficienten von (10), d. h. rational durch $a_i,\alpha_i,y$ ausdrücken, vorausgesetzt, dass nur eine solche gemeinsame Wurzel vorhanden ist.

Dies wird dann eintreten, wenn die $n$ Werthe $y_1,y_2,\ldots,y_n$ von einander verschieden sind. Dann erhält man einen Ausdruck
\[
x=\beta_0+\beta_1y+\beta_2y^2+\cdots+\beta_{n-1}y^{n-1},
\]
der gültig ist, wenn für $x$ einer der Werthe (6) und für $y$ der zugehörige Werth (7) gesetzt wird. Die Grössen $\beta_i$ sind aus den $a_i$ und $\alpha_i$ rational zusammengesetzt.

Kommen aber unter den Werthen $y_1,y_2,\ldots,y_n$ mehrere gleiche vor, so erfordert die Bestimmung der zugehörigen Werthe der $x$ noch die Auflösung einer Gleichung des entsprechenden Grades. Dieser Grad ist aber immer kleiner als $n$, da die Function $\varphi(x)$ höchstens für $n-1$ Werthe von $x$ denselben Werth $y$ erhalten kann. Wenn also die Gleichung
\begin{equation}
\Phi(y)=0
\tag{11}
\end{equation}
vollständig gelöst ist, so ist damit auch die Gleichung (1) vollständig gelöst [wenn (11) $n$ verschiedene Wurzeln hat] oder wenigstens auf die Lösung einer Gleichung niederen Grades zurückgeführt [wenn (11) gleiche Wurzeln hat].

Demnach betrachten wir (11) als eine Umformung der Gleichung (1), und die Bildung von (11) aus (1) heisst die Tschirnhausen-Transformation der Gleichung (1).

Der Zweck einer solchen Transformation ist, die Coefficienten $\alpha_i$ so zu bestimmen, dass $\Phi(y)$ eine einfachere Gestalt erhält als $f(x)$.

Man kann die Coefficienten $b_\nu$ bilden, wenn man die Potenzsummen der $y_i$ kennt (\S\,42).

Bezeichnen wir die Summen der $m$ten Potenzen der $x_i$ mit $s_m$, der $y_i$ mit $\sigma_m$, also
\[
s_m=Sx_i^m,
\qquad
\sigma_m=Sy_i^m,
\]
so können wir $\sigma_m$ durch die $s_m$ auf folgende Art bilden. Wir entwickeln zunächst $y^m$ nach dem polynomischen Lehrsatz (\S\,12), setzen also
\begin{equation}
y^m=\sum \frac{\Pi(m)}{\Pi(\nu_0)\Pi(\nu_1)\cdots\Pi(\nu_{n-1})}
\alpha_0^{\nu_0}\alpha_1^{\nu_1}\cdots\alpha_{n-1}^{\nu_{n-1}}
 x^{\nu_1+2\nu_2+\cdots+(n-1)\nu_{n-1}},
\tag{12}
\end{equation}
worin die Summe sich auf alle nicht negativen ganzzahligen Werthe von $\nu_0,\nu_1,\ldots,\nu_{n-1}$ erstreckt, die der Bedingung
\begin{equation}
\nu_0+\nu_1+\cdots+\nu_{n-1}=m
\tag{13}
\end{equation}
genügen. Setzen wir dann in (12) $x_1,x_2,
\ldots,x_n$ für $x$ und bilden die Summe der so erhaltenen Werthe von $y_1^m,y_2^m,\ldots,y_n^m$, so folgt
\begin{equation}
\sigma_m=\sum \frac{\Pi(m)}{\Pi(\nu_0)\Pi(\nu_1)\cdots\Pi(\nu_{n-1})}
 s_{\nu_1+2\nu_2+\cdots+(n-1)\nu_{n-1}}
 \alpha_0^{\nu_0}\alpha_1^{\nu_1}\cdots\alpha_{n-1}^{\nu_{n-1}}.
\tag{14}
\end{equation}
Hierdurch ist $\sigma_m$ als ganze homogene Function $m$ter Ordnung der $\alpha_0,\alpha_1,\ldots,\alpha_{n-1}$ dargestellt. Eine etwas vereinfachte Schreibweise erhält man nach der zweiten Darstellung von Formen (\S\,15). Lassen wir $\mu_1,\mu_2,\ldots,\mu_m$ unabhängig von einander die Reihe der Zahlen $0,1,\ldots,n-1$ durchlaufen, so wird, wenn in einer Combination $\nu_0$ mal der Index $0$, $\nu_1$ mal der Index $1$, $\nu_2$ mal der Index $2$, $\ldots$, $\nu_{n-1}$ mal der Index $n-1$ vorkommt,
\[
\mu_1+\mu_2+\cdots+\mu_m=\nu_1+2\nu_2+\cdots+(n-1)\nu_{n-1},
\]
und danach kann der Ausdruck für $\sigma_m$ so dargestellt werden:
\begin{equation}
\sigma_m=\sum_{\mu} s_{\mu_1+\mu_2+\cdots+\mu_m}\alpha_{\mu_1}\alpha_{\mu_2}\cdots\alpha_{\mu_m}.
\tag{15}
\end{equation}
Beispielsweise wird
\begin{equation}
\begin{aligned}
\sigma_1&=\alpha_0s_0+\alpha_1s_1+\cdots+\alpha_{n-1}s_{n-1},\\
\sigma_2&=\sum_{i,k}s_{i+k}\alpha_i\alpha_k,\\
\sigma_3&=\sum_{h,i,k}s_{h+i+k}\alpha_h\alpha_i\alpha_k,
\end{aligned}
\tag{16}
\end{equation}
worin $h,i,k$ die Reihe der Zahlen $0,1,2,\ldots,n-1$ durchlaufen.

\sect{53}{Anwendung auf die cubischen und biquadratischen Gleichungen}

Das Ziel, das schon Tschirnhausen bei dieser Transformation im Auge gehabt hat, bestand darin, durch Bestimmung der Substitutionscoefficienten $\alpha_0,\alpha_1,\ldots,\alpha_{n-1}$ in der umgeformten Gleichung $\Phi(y)=0$ so viele Glieder zum Verschwinden zu bringen, dass die Gleichung lösbar wird. Es gelingt dies leicht bei den Gleichungen dritten und vierten Grades, wenn es auch nicht einfach ist, und nicht ohne weitläufige Rechnungen möglich scheint, die sonst bekannten Formeln auf diesem Wege abzuleiten.

Nehmen wir zunächst die cubische Gleichung
\begin{equation}
x^3+ax^2+bx+c=0,
\tag{1}
\end{equation}
die wir durch die Substitution
\begin{equation}
y=\alpha+\beta x+\gamma x^2
\tag{2}
\end{equation}
umformen. Um die Gleichung für $y$
\begin{equation}
y^3+Ay^2+By+C=0
\tag{3}
\end{equation}
auf eine reine cubische Gleichung zurückzuführen, haben wir die Verhältnisse $\alpha:\beta:\gamma$ aus den beiden Gleichungen
\begin{equation}
A=0,
\qquad
B=0
\tag{4}
\end{equation}
zu bestimmen, von denen die erste linear, die zweite quadratisch ist.

Die Gleichungen (4) sind gleichbedeutend mit
\begin{equation}
\sigma_1=0,
\qquad
\sigma_2=0
\tag{5}
\end{equation}
und ergeben also nach (16) des vorigen Paragraphen
\[
\alpha s_0+\beta s_1+\gamma s_2=0,
\]
\[
\alpha^2s_0+2\alpha\beta s_1+2\alpha\gamma s_2+\beta^2s_2+2\beta\gamma s_3+\gamma^2s_4=0,
\]
worin $s_0=3$ ist. Multiplicirt man die letztere Gleichung mit $s_0$ und zieht das Quadrat der ersten davon ab, so folgt
\begin{equation}
\beta^2(s_0s_2-s_1^2)+2\beta\gamma(s_0s_3-s_1s_2)+\gamma^2(s_0s_4-s_2^2)=0.
\tag{6}
\end{equation}
Die Discriminante dieser quadratischen Gleichung
\[
(s_0s_3-s_1s_2)^2-(s_0s_2-s_1^2)(s_0s_4-s_2^2)
\]
giebt entwickelt
\[
-s_0(s_0s_2s_4-s_0s_3^2-s_1^2s_4-s_2^3+2s_1s_2s_3),
\]
ist also, wenn $D$ die Discriminante der cubischen Gleichung (1) bedeutet [\S\,46, (9)], gleich $-3D$, und die Auflösung von (6) giebt also
\begin{equation}
\frac{\beta}{\gamma}=\frac{-(s_0s_3-s_1s_2)+\sqrt{-3D}}{s_0s_2-s_1^2}.
\tag{7}
\end{equation}
Damit ist die Gleichung (3) auf eine reine Gleichung zurückgeführt, die zu ihrer Lösung nur noch das Ziehen einer Cubikwurzel erfordert. Welches Vorzeichen wir der in (7) vorkommenden Quadratwurzel geben wollen, steht in unserem Belieben.

Um die Tschirnhausen-Transformation für die biquadratische Gleichung zu benutzen, kann man etwa so verfahren, dass man in der umgeformten Gleichung
\begin{equation}
y^4+b_1y^3+b_2y^2+b_3y+b_4=0
\tag{8}
\end{equation}
die $\alpha_0,\alpha_1,\alpha_2,\alpha_3$ so bestimmt, dass
\begin{equation}
b_1=0,
\qquad
b_3=0
\tag{9}
\end{equation}
wird, wodurch (8) in eine quadratische Gleichung für $y^2$ übergeht, aus deren Wurzeln noch die Quadratwurzeln gezogen werden müssen, so dass, nachdem die Gleichungen (9) gelöst sind, noch zweimal eine Quadratwurzel gezogen werden muss.

Wenn man mittelst der ersten Gleichung (9) aus der zweiten $\alpha_0$ eliminirt, so entsteht eine homogene cubische Gleichung zwischen $\alpha_1,\alpha_2,\alpha_3$. Man kann daher eines von den beiden Verhältnissen $\alpha_1:\alpha_2:\alpha_3$ beliebig annehmen, z. B. $\alpha_3=0$ setzen, und erhält zur Bestimmung des anderen eine cubische Gleichung. Dadurch ist dann die Lösung der biquadratischen Gleichung auf die einer cubischen und auf Quadratwurzeln zurückgeführt.

\sect{54}{Die Tschirnhausen-Transformation der Gleichung fünften Grades}

Will man auf die Gleichung fünften Grades
\begin{equation}
x^5+a_1x^4+a_2x^3+a_3x^2+a_4x+a_5=0
\tag{1}
\end{equation}
die Tschirnhausen-Transformation anwenden, so hat man
\begin{equation}
y=\alpha_0+\alpha_1x+\alpha_2x^2+\alpha_3x^3+\alpha_4x^4
\tag{2}
\end{equation}
zu setzen, und erhält die Transformirte
\begin{equation}
y^5+b_1y^4+b_2y^3+b_3y^2+b_4y+b_5=0.
\tag{3}
\end{equation}
Der nächstliegende Gedanke wäre nun, die Verhältnisse der fünf Unbekannten $\alpha$ so zu bestimmen, dass
\begin{equation}
b_1=0,
\qquad b_2=0,
\qquad b_3=0,
\qquad b_4=0
\tag{4}
\end{equation}
würde, wodurch (3) auf eine reine Gleichung zurück käme, und dies war wohl auch das Ziel, das Tschirnhausen im Auge hatte. Aber von den vier Gleichungen (4) ist die erste linear, die zweite quadratisch, die dritte vom dritten und die vierte vom vierten Grade. Die Endgleichung, die man daraus ableiten kann, wird daher (nach \S\,50) vom Grade $2\cdot3\cdot4=24$, und daher kann daraus für die Lösung der Gleichung fünften Grades kein Nutzen gezogen werden. Man sucht daher zunächst nur den beiden Gleichungen zu genügen
\begin{equation}
b_1=0,
\qquad b_2=0,
\tag{5}
\end{equation}
wodurch die Gleichung (3) in eine Form übergeht, die nach F. Klein (Vorlesungen über das Ikosaeder) eine Hauptgleichung fünften Grades genannt wird.

Zur Befriedigung der beiden Gleichungen (5) haben wir die Verfügung über die vier Verhältnisse $\alpha_1:\alpha_2:\alpha_3:\alpha_4$, und wir können diese Grössen noch auf mannigfaltige Art zur Vereinfachung der Gleichung fünften Grades benutzen.

Wir kommen in einem späteren Abschnitt ausführlicher auf diesen Gegenstand zurück und beschränken uns hier darauf, die Ziele im Allgemeinen zu bezeichnen. Der Anschaulichkeit wegen bedienen wir uns einer geometrischen Einkleidung, die übrigens zum Verständniss der algebraischen Theorie, wie sie später gegeben werden soll, durchaus nicht wesentlich ist.

Wenn wir mit Hülfe der linearen Gleichung $b_1=0$ von den fünf Grössen $\alpha$ die eine, etwa $\alpha_0$, durch die übrigen ausdrücken, so können wir die vier übrigen $\alpha_1,\alpha_2,\alpha_3,\alpha_4$ als homogene Coordinaten eines Raumpunktes betrachten. Die Gleichung $b_2=0$ stellt dann eine Fläche zweiter Ordnung dar, $b_3=0$ eine Fläche dritter Ordnung, $b_4=0$ eine Fläche vierter Ordnung, $b_5=0$ eine Fläche fünfter Ordnung.

Wollte man also $b_2,b_3,b_4$ zugleich zu Null machen und dadurch die gegebene Gleichung auf eine reine reduciren, so müsste man einen der 24 Schnittpunkte von drei Flächen zweiter, dritter, vierter Ordnung bestimmen, der von einer Gleichung $24$ten Grades abhängt.

Statt nun so zu verfahren, sucht man auf der Fläche $b_2=0$ eine gerade Linie zu bestimmen. Solcher gerader Linien giebt es auf jeder Fläche zweiter Ordnung eine oder zwei Schaaren (reell oder imaginär), und man kann eine dieser Geraden durch Quadratwurzeln bestimmen, wie weiter unten ausgeführt werden soll. Die Ermittelung eines Schnittpunktes einer dieser geraden Linien mit der Fläche $b_3=0$ führt dann auf eine Gleichung dritten Grades, und wenn man noch durch Bestimmung eines gemeinsamen Factors der $\alpha$ den Coefficienten $b_4=1$ macht, so erhält die Gleichung für $y$ die Gestalt
\begin{equation}
y^5+y+b=0,
\tag{6}
\end{equation}
wodurch $y$ als Function einer Variabeln $b$ bestimmt ist.

Diese Gleichungsform wird gewöhnlich nach dem englischen Mathematiker Jerrard genannt, der sie im Jahre 1834 bekannt machte. Sie ist aber schon viel früher (1786) von E. J. Bring in Lund publicirt worden.\footnote{Vgl. F. Klein, Vorlesungen über das Ikosaëder, S. 143.} Wir nennen sie daher die Bring-Jerrard'sche Form.

Ebensogut könnte man durch eine Gleichung vierten Grades $b_4=0$ machen, und würde eine zweite Normalform der Gleichung fünften Grades erhalten:
\begin{equation}
y^5+y^2+b=0,
\tag{7}
\end{equation}
die die gleiche Berechtigung hat, wie die Bring-Jerrard'sche Form.

Für die Lösung der Gleichung fünften Grades freilich wird durch diese Betrachtungen nichts gewonnen. Ihr Nutzen besteht darin, dass die Gleichung fünften Grades auf eine Normalform zurückgeführt wird, die nur von einem veränderlichen Coefficienten abhängt. Dies kann auf unendlich viele verschiedene Arten erreicht werden und gelingt am einfachsten, nämlich ohne Lösung einer Gleichung dritten oder vierten Grades, wenn man die Gleichung fünften Grades $b_5=0$ selbst als eine Umformung der gegebenen Gleichung fünften Grades betrachtet. Hat man nämlich die Gleichung $b_5=0$ gelöst, so wird einer der fünf Werthe von $y$ gleich $0$, und die gegebene Gleichung fünften Grades ist damit zugleich gelöst.

\end{document}
